\markboth{INTRODUCTION}{INTRODUCTION}
\section{Introduction}\label{sec:intro}

Wearable devices from Garmin, Polar, and Apple now capture heart rate, GPS trajectories, sleep stages, heart-rate variability (HRV), and stress levels continuously~\citep{seckin_review_2023}. Strava alone aggregates activity data from over 180~million users~\citep{strava_trends_2025}. The result is an unprecedented volume of training data---yet it remains fragmented across proprietary ecosystems with incompatible interfaces, making holistic analysis impractical for anyone without a sports-science background~\citep{bourdon_monitoring_2017}.

A concrete example illustrates this. An endurance cyclist we encountered during user research exports his Strava data after every ride, uploads it to ChatGPT, and asks the model to analyse his progress. The process is both slow and incomplete: the general-purpose LLM has no access to his Garmin recovery data, the weather forecast, or his calendar. More broadly, deciding whether to train on a given day requires checking a Garmin watch for HRV and sleep quality, Strava for cumulative load, a weather service for conditions, and a calendar for scheduling. Each tool requires a separate application, and the integration is left entirely to the athlete.

Large language models have demonstrated the ability to reason across heterogeneous information sources when equipped with tool-calling capabilities~\citep{schick_toolformer_2023, qu_tool_2024}. The emergence of open interoperability protocols---Anthropic's Model Context Protocol (MCP)~\citep{anthropic_mcp_2024} for tool discovery and invocation, and Google's Agent-to-Agent (A2A)~\citep{google_a2a_2025} for inter-agent coordination---makes it feasible to build a system that connects these fragmented sources behind a single conversational interface.

\subsection{Problem Statement and Approach}\label{sec:intro:problem}

This paper presents \textit{Training Copilot}, an open-source sports analytics platform that addresses three problems:

\textbf{Data fragmentation.} Training and wellness data is distributed across platforms with incompatible authentication models, data formats, and access patterns. Sports science consensus statements emphasise that effective load monitoring requires combining external and internal load from multiple sources~\citep{bourdon_monitoring_2017, halson_monitoring_2014}, yet the consumer tools that collect this data do not provide a unified view. Training Copilot solves this with MCP as a uniform data layer: each external API is encapsulated in an independent MCP server, and a single host client discovers and routes tool calls without hardcoding provider-specific logic.

\textbf{Lack of cross-domain reasoning.} The platforms we examined---Strava, Garmin Connect, Intervals.icu, TrainingPeaks---display metrics in isolation: a sleep score here, a training load chart there. An athlete's readiness to train depends on the interplay of recovery status, cumulative workload, environmental conditions, and personal schedule~\citep{bourdon_monitoring_2017, gabbett_training_2016}, but none of these tools synthesise across all four dimensions. Training Copilot addresses this with a multi-agent system: five LangGraph~\citep{langchain_langgraph_2024} ReAct~\citep{yao_react_2023} agents, each scoped to a specific domain (recovery, load, context, routes, fitness knowledge), coordinated by an orchestrator via the A2A protocol.

\textbf{Rigid architectures.} The platforms listed above are tightly coupled to their data sources---adding a new source (e.g., a nutrition tracker) would require code changes throughout their stack. Training Copilot's MCP-based design reduces adding a new data source to one server file and one configuration line, consistent with microservice design principles~\citep{dragoni_microservices_2017}.

\subsection{Contributions}\label{sec:intro:contributions}

\begin{itemize}
\item An MCP-based integration architecture that achieves single-line extensibility for new data sources, verified empirically during development by adding the Telegram bridge and flythrough servers without changes to the host, agents, or UI.
\item A multi-agent system with five domain-scoped specialists coordinated via A2A, where scoping each specialist to at most two MCP servers addresses the tool-selection degradation observed with large tool sets~\citep{qu_tool_2024}.
\item A RAG-based fitness knowledge agent that grounds exercise-science advice in a local vector database of public-domain literature, with explicit citation preservation through a post-processing step.
\item An automated evaluation framework using MLflow conversation simulation~\citep{mlflow_multiturn_2025} with ten persona-driven test cases and five scoring dimensions, including a deterministic trace-based grounding scorer.
\end{itemize}

\subsection{Structure of This Paper}\label{sec:intro:structure}

Section~\ref{sec:market} analyses the sports analytics market and identifies the gap Training Copilot addresses. Section~\ref{sec:sota} reviews the state of the art in agentic AI, tool integration protocols, sports analytics, and retrieval-augmented generation. Section~\ref{sec:architecture} presents the system architecture. Section~\ref{sec:data} describes the data sources and their integration. Section~\ref{sec:agents} details the agent interaction patterns. Section~\ref{sec:evaluation} describes the evaluation methodology. Section~\ref{sec:discussion} discusses trade-offs, limitations, and ethical considerations. Section~\ref{sec:conclusion} concludes with a summary and directions for future work.
